\documentclass[10pt,a4paper]{amsart}
\usepackage{amsmath,amssymb,amsthm,graphicx}
% Retained preamble for the Stacks Project chapter snapshot in this directory.
% The chapter's own preamble.tex is a copy of the shared Stacks Project preamble; the block
% below is sliced MECHANICALLY, line for line, from it. Nothing was added, removed or
% redefined; the wrapper class, the xr-hyper/hyperref wiring and the externaldocument list
% are omitted because the assembler supplies the document class.
%
% lines taken from preamble.tex: 3, 5, 9, 11, 12, 156, 157, 158, ... (51 lines total)
%
% package coverage restored from the chapter's own preamble (2026-10-13): xy, xy-2cell

\usepackage[all]{xy}
\xyoption{2cell}
\UseAllTwocells

\newenvironment{reference}{\comment}{\endcomment}
\newenvironment{slogan}{\comment}{\endcomment}
\newenvironment{history}{\comment}{\endcomment}
\theoremstyle{plain}
\newtheorem{theorem}[subsection]{Theorem}
\newtheorem{proposition}[subsection]{Proposition}
\newtheorem{lemma}[subsection]{Lemma}
\theoremstyle{definition}
\newtheorem{definition}[subsection]{Definition}
\newtheorem{example}[subsection]{Example}
\newtheorem{exercise}[subsection]{Exercise}
\newtheorem{situation}[subsection]{Situation}
\theoremstyle{remark}
\newtheorem{remark}[subsection]{Remark}
\newtheorem{remarks}[subsection]{Remarks}
\numberwithin{equation}{subsection}
\def\lim{\mathop{\mathrm{lim}}\nolimits}
\def\colim{\mathop{\mathrm{colim}}\nolimits}
\def\Spec{\mathop{\mathrm{Spec}}}
\def\Hom{\mathop{\mathrm{Hom}}\nolimits}
\def\Ext{\mathop{\mathrm{Ext}}\nolimits}
\def\SheafHom{\mathop{\mathcal{H}\!\mathit{om}}\nolimits}
\def\SheafExt{\mathop{\mathcal{E}\!\mathit{xt}}\nolimits}
\def\Sch{\mathit{Sch}}
\def\Mor{\mathop{\mathrm{Mor}}\nolimits}
\def\Ob{\mathop{\mathrm{Ob}}\nolimits}
\def\Sh{\mathop{\mathit{Sh}}\nolimits}
\def\NL{\mathop{N\!L}\nolimits}
\def\CH{\mathop{\mathrm{CH}}\nolimits}
\def\proetale{{pro\text{-}\acute{e}tale}}
\def\etale{{\acute{e}tale}}
\def\QCoh{\mathit{QCoh}}
\def\Ker{\mathop{\mathrm{Ker}}}
\def\Im{\mathop{\mathrm{Im}}}
\def\Coker{\mathop{\mathrm{Coker}}}
\def\Coim{\mathop{\mathrm{Coim}}}
\DeclareMathSymbol{\boxtimes}{\mathbin}{AMSa}{"02}
\def\QCohstack{\mathcal{QC}\!\mathit{oh}}
\def\Cohstack{\mathcal{C}\!\mathit{oh}}
\def\Spacesstack{\mathcal{S}\!\mathit{paces}}
\def\Quotfunctor{\mathrm{Quot}}
\def\Hilbfunctor{\mathrm{Hilb}}
\def\Curvesstack{\mathcal{C}\!\mathit{urves}}
\def\Polarizedstack{\mathcal{P}\!\mathit{olarized}}
\def\Complexesstack{\mathcal{C}\!\mathit{omplexes}}
\def\Pic{\mathop{\mathrm{Pic}}\nolimits}
\def\Picardstack{\mathcal{P}\!\mathit{ic}}
\def\Picardfunctor{\mathrm{Pic}}
\def\Deformationcategory{\mathcal{D}\!\mathit{ef}}
\usepackage{cleveref}
\crefname{originallemma}{original lemma}{original lemmas}
\Crefname{originallemma}{Original lemma}{Original lemmas}
\makeatletter
\newlabel{algebra-lemma-characterize-zero-local}{{00HN}{source}{}{sourcepointer-1}{}}
\newlabel{algebra-lemma-characterize-zero-local@cref}{{[originallemma][0][]00HN}{[1][0][]source}{}{}{}}
\makeatother
\crefname{originallemma}{original lemma}{original lemmas}
\Crefname{originallemma}{Original lemma}{Original lemmas}
\makeatletter
\newlabel{algebra-lemma-etale-makes-quasi-finite-finite-one-prime}{{00UJ}{source}{}{sourcepointer-2}{}}
\newlabel{algebra-lemma-etale-makes-quasi-finite-finite-one-prime@cref}{{[originallemma][0][]00UJ}{[1][0][]source}{}{}{}}
\makeatother
\crefname{originallemma}{original lemma}{original lemmas}
\Crefname{originallemma}{Original lemma}{Original lemmas}
\makeatletter
\newlabel{lemma-dualizing-flat-base-change}{{0BT3}{source}{}{sourcepointer-3}{}}
\newlabel{lemma-dualizing-flat-base-change@cref}{{[originallemma][0][]0BT3}{[1][0][]source}{}{}{}}
\makeatother
\crefname{originallemma}{original lemma}{original lemmas}
\Crefname{originallemma}{Original lemma}{Original lemmas}
\makeatletter
\newlabel{lemma-dualizing-product}{{0BT5}{source}{}{sourcepointer-4}{}}
\newlabel{lemma-dualizing-product@cref}{{[originallemma][0][]0BT5}{[1][0][]source}{}{}{}}
\makeatother
\crefname{originallemma}{original lemma}{original lemmas}
\Crefname{originallemma}{Original lemma}{Original lemmas}
\makeatletter
\newlabel{lemma-noether-different-base-change}{{0BVM}{source}{}{sourcepointer-5}{}}
\newlabel{lemma-noether-different-base-change@cref}{{[originallemma][0][]0BVM}{[1][0][]source}{}{}{}}
\makeatother
\crefname{originallemma}{original lemma}{original lemmas}
\Crefname{originallemma}{Original lemma}{Original lemmas}
\makeatletter
\newlabel{lemma-noether-different-product}{{0BVL}{source}{}{sourcepointer-6}{}}
\newlabel{lemma-noether-different-product@cref}{{[originallemma][0][]0BVL}{[1][0][]source}{}{}{}}
\makeatother
\crefname{originallemma}{original lemma}{original lemmas}
\Crefname{originallemma}{Original lemma}{Original lemmas}
\makeatletter
\newlabel{lemma-noether-pairing-compatibilities}{{0BVR}{source}{}{sourcepointer-7}{}}
\newlabel{lemma-noether-pairing-compatibilities@cref}{{[originallemma][0][]0BVR}{[1][0][]source}{}{}{}}
\makeatother
\crefname{originallemma}{original lemma}{original lemmas}
\Crefname{originallemma}{Original lemma}{Original lemmas}
\makeatletter
\newlabel{remark-construction-pairing}{{0BVP}{source}{}{sourcepointer-8}{}}
\newlabel{remark-construction-pairing@cref}{{[originallemma][0][]0BVP}{[1][0][]source}{}{}{}}
\makeatother
\title{Stacks, Discriminants: lemma-noether-pairing-flat-quasi-finite}
\author{The Stacks Project Authors}
\date{stacks-project:discriminant.tex:local-snapshot-20261012; GNU FDL 1.2 or later}
\begin{document}
\maketitle
\noindent\textit{Editorial note.} Statement span is lines 1427-1432 of sources/stacks-discriminant/pinned-original-source.tex; proof span is lines 1434-1482, the complete author proof (47 body lines).\par\medskip
\noindent\textit{Editorial note.} Navigation stubs: 8 cited-but-not-included labels are supplied as non-load-bearing pointers with their pinned tags where the snapshot tag index has them.\par\medskip
\noindent\textit{Editorial note.} extra\_packages is empty on purpose: the retained preamble already loads xy with options and a second declaration clashes in the assembler wrapper.\par\medskip
\noindent\textit{Editorial note.} Text is retained verbatim, including the author's own typos and placeholders.\par\medskip
\noindent\textit{Editorial note.} Queue index 7 of queue\_discriminant.json (pb=2114, stmt 1427-1432, proof 1434-1482).\par\medskip
\noindent\textit{Allowed standard prerequisites (editorial).}\begin{itemize}
\item Explicit prerequisite boundary. The cited inputs are prerequisites; the argument in the retained span is not replaced by any cited result.
\end{itemize}\medskip
\ifdefined\hypertarget\hypertarget{sourcepointer-1}{}\fi
\noindent\textit{Source reference.} Original lemma 00HN: sources/stacks-discriminant/pinned-original-source.tex, source lines 1482--1482. Cited by the retained statement or proof but not part of the collected spans; supplied only as a navigation pointer. The argument does not depend on its content.\par\medskip
\ifdefined\hypertarget\hypertarget{sourcepointer-2}{}\fi
\noindent\textit{Source reference.} Original lemma 00UJ: sources/stacks-discriminant/pinned-original-source.tex, source lines 1482--1482. Cited by the retained statement or proof but not part of the collected spans; supplied only as a navigation pointer. The argument does not depend on its content.\par\medskip
\ifdefined\hypertarget\hypertarget{sourcepointer-3}{}\fi
\noindent\textit{Source reference.} Original lemma 0BT3: sources/stacks-discriminant/pinned-original-source.tex, source lines 450--454. Cited by the retained statement or proof but not part of the collected spans; supplied only as a navigation pointer. The argument does not depend on its content.\par\medskip
\ifdefined\hypertarget\hypertarget{sourcepointer-4}{}\fi
\noindent\textit{Source reference.} Original lemma 0BT5: sources/stacks-discriminant/pinned-original-source.tex, source lines 493--497. Cited by the retained statement or proof but not part of the collected spans; supplied only as a navigation pointer. The argument does not depend on its content.\par\medskip
\ifdefined\hypertarget\hypertarget{sourcepointer-5}{}\fi
\noindent\textit{Source reference.} Original lemma 0BVM: sources/stacks-discriminant/pinned-original-source.tex, source lines 1311--1320. Cited by the retained statement or proof but not part of the collected spans; supplied only as a navigation pointer. The argument does not depend on its content.\par\medskip
\ifdefined\hypertarget\hypertarget{sourcepointer-6}{}\fi
\noindent\textit{Source reference.} Original lemma 0BVL: sources/stacks-discriminant/pinned-original-source.tex, source lines 1295--1303. Cited by the retained statement or proof but not part of the collected spans; supplied only as a navigation pointer. The argument does not depend on its content.\par\medskip
\ifdefined\hypertarget\hypertarget{sourcepointer-7}{}\fi
\noindent\textit{Source reference.} Original lemma 0BVR: sources/stacks-discriminant/pinned-original-source.tex, source lines 1386--1411. Cited by the retained statement or proof but not part of the collected spans; supplied only as a navigation pointer. The argument does not depend on its content.\par\medskip
\ifdefined\hypertarget\hypertarget{sourcepointer-8}{}\fi
\noindent\textit{Source reference.} Original lemma 0BVP: sources/stacks-discriminant/pinned-original-source.tex, source lines 1362--1369. Cited by the retained statement or proof but not part of the collected spans; supplied only as a navigation pointer. The argument does not depend on its content.\par\medskip
\begin{lemma}
\label{lemma-noether-pairing-flat-quasi-finite}
Let $A \to B$ be a flat quasi-finite homomorphism of Noetherian rings.
The pairing of Remark \ref{remark-construction-pairing} induces an isomorphism
$J \to \Hom_B(\omega_{B/A}, B)$.
\end{lemma}

\begin{proof}
We first prove this when $A \to B$ is finite and flat. In this case we can
localize on $A$ and assume $B$ is finite free as an $A$-module. Let
$b_1, \ldots, b_n$ be a basis of $B$ as an $A$-module and denote
$b_1^\vee, \ldots, b_n^\vee$ the dual basis of $\omega_{B/A}$. Note that
$\sum b_i \otimes c_i \in J$ maps to the element of $\Hom_B(\omega_{B/A}, B)$
which sends $b_i^\vee$ to $c_i$. Suppose $\varphi : \omega_{B/A} \to B$
is $B$-linear. Then we claim that $\xi = \sum b_i \otimes \varphi(b_i^\vee)$
is an element of $J$. Namely, the $B$-linearity of $\varphi$
exactly implies that $(b \otimes 1)\xi = (1 \otimes b)\xi$ for all $b \in B$.
Thus our map has an inverse and it is an isomorphism.

\medskip\noindent
Let $\mathfrak q \subset B$ be a prime lying over $\mathfrak p \subset A$.
We will show that the localization
$$
J_\mathfrak q
\longrightarrow
\Hom_B(\omega_B/A, B)_\mathfrak q
$$
is an isomorphism.
This suffices by Algebra, Lemma \ref{algebra-lemma-characterize-zero-local}.
By
Algebra, Lemma \ref{algebra-lemma-etale-makes-quasi-finite-finite-one-prime}
we can find an \'etale ring map $A \to A'$ and a prime
ideal $\mathfrak p' \subset A'$ lying over $\mathfrak p$
such that $\kappa(\mathfrak p') = \kappa(\mathfrak p)$ and
such that
$$
B' = B \otimes_A A' = C \times D
$$
with $A' \to C$ finite and such that the unique prime $\mathfrak q'$
of $B \otimes_A A'$ lying over $\mathfrak q$ and $\mathfrak p'$
corresponds to a prime of $C$. Let $J'$ be the annihilator of
$\Ker(B' \otimes_{A'} B' \to B')$. By
Lemmas \ref{lemma-dualizing-flat-base-change},
\ref{lemma-noether-different-base-change}, and
\ref{lemma-noether-pairing-compatibilities}
the map $J' \to \Hom_{B'}(\omega_{B'/A'}, B')$
is gotten by applying the functor $- \otimes_B B'$
to the map $J \to \Hom_B(\omega_{B/A}, B)$.
Since $B_\mathfrak q \to B'_{\mathfrak q'}$ is faithfully flat
it suffices to prove the result for $(A' \to B', \mathfrak q')$.
By Lemmas \ref{lemma-dualizing-product},
\ref{lemma-noether-different-product}, and
\ref{lemma-noether-pairing-compatibilities}
this reduces us to the case proved in the first
paragraph of the proof.
\end{proof}

\end{document}
